% Folder 134-2, batch 09 — pages 161–180 (position in the folder; see below).
% Pass: Opus 5.5 (claude-opus-5-5), 2026-10-15 — first pass, unchecked against the pages by a human.
%
% Special rules for folder 134-2 (issues #44 and #57, the latter as decided by the
% user on 2026-10-15) applied: his own letters not in print are transcribed in full.
%
% Numbering: \page{} numbers the position in the folder, as in batches 6 and 7. From
% p. 115 on the archivists' pencil runs one ahead of the position: here p. 161 is
% pencilled 162, p. 180 is pencilled 181. His own page numbers stand in the top margin
% and are recorded once per page in a \note.
%
% Contents of the batch:
%  - pp. 161-180 (his pp. 17 to 36; « 17 » itself is not visible on p. 161): the
%    HANDWRITTEN FRENCH ORIGINAL of the letter to Lawrence Breen, « Villecun 17/19
%    juillet 1975 » (« Cher Larry »), begun in batch 8 (p. 145) and continued in
%    batch 10 (the letter is signed on p. 184). Transcribed in full. The batch opens
%    mid-sentence, inside the bracket on « Lang's trick » of section (D), and runs
%    through (E) local class field theory, (F) the fppf topology, the answer to
%    Breen's question 7 (homotopy type of a topos, locally constant n-stacks,
%    stratifications, tubular neighbourhoods, tame topology) and stops mid-sentence
%    in the « fine homotopy type » paragraph (English version: batch 4, p. 75).
%  - No unrelated verso, no blank page, no page absent.
%  - The English typed version (batches 3-4, pp. 56-77) was used only to follow the
%    argument and to place some interlinear insertions; where it settled an
%    \uncertain{} or the order of insertions, the note says so. No \ill{} was filled
%    from it. Differences between the French and the typed English (e.g. « n-systèmes
%    locaux », « schéma en anneaux sur k défini par V », the subscripts of
%    Δ_0 RΓ_K(F) ≃ RΓ_K(DF[2])) are kept as on the French page and noted.
%  - The English marginal calls « (App. 12) » to « (App. 18) » and footnote calls
%    (28)-(32) of the typescript are NOT on the French original.
%
% Letters in this batch: Grothendieck to L. Breen, 17/19.7.1975, French autograph,
% his pp. 17-36 (transcribed).
\documentclass[11pt,a4paper]{article}
\input{../preamble/grothendieck}

\folder{134-2}
\batch{9}
\pages{161}{180}
\foldertitle{[Chapitre I : Take off (pages 1 à 65) et table des matières provisoire] : tapuscrits annotés (19/02-22/02/1983), note manuscrite (s.d.), copies de lettre (1975, s.d.).}
\dating{1975-[1983]}
\watermark{Édition de démonstration}

\begin{document}

\section{Lettre à L. Breen, Villecun 17/19 juillet 1975 (original manuscrit, suite)}

\page{161}
\note{p.~17 de l'auteur, d'après la suite de sa pagination (le numéro n'est pas visible en tête de la page) ; crayon des archivistes : 162. La page commence au milieu d'une phrase, dans le crochet sur « l'astuce de Lang » ouvert au lot~8.}
stant -- $\pi_1^{\mathrm{ar}}(J) \simeq H^0(k, J)$] sous une forme cohomologique, en énonçant que l'on a
\[
\Delta_0 \mathbb{R}\Gamma_k(J^*) \simeq \mathbb{R}\Gamma_k(\Delta J^*[1])
\]
où $\Delta$ est la dualité de Serre, $\Delta_0$ la dualité de Pontrjagin des groupes abéliens top. totalement discontinus (dualité à valeurs dans $\mathbb{Q}/\mathbb{Z}$), $J^*$ un complexe de ind-pro-groupes algébriques sur $k$. Tenant compte de cette « formule de dualité de Lang » et appliquant $\mathbb{R}\Gamma_k$ à la formule de dualité \add{des corps de classes} géométrique, on trouve la « formule de dualité des corps de classes arithmétique »
\[
\boxed{\Delta_0(H_!(X, F)) \simeq H^*(X, D(F)[3])}
\]
(isom. de groupes topologiques tot. discontinus).
\note{au-dessus de $H^0(k, J)$, deux lettres (« $\vee$ », « $+$ » ?) dont le rôle n'est pas clair. « $\mathbb{R}\Gamma_k$ » : l'indice du second membre ressemble à un $h$. La formule finale est encadrée. Le sens de « sous une forme cohomologique » (l'astuce de Lang « dans sa forme cohomologique ») est confirmé par la version anglaise, lot~4, p.~64.}

Autre remarque : lorsque $F$ n'est pas un « faisceau étale », mais a une structure continue comme $\alpha_p$, il faut faire attention dans la définition de $Rf_!(F)$, pour $X$ non complet, à partir de la compactifiée \add{$\tilde{X}$} ; ainsi, si $F$ provient d'un faisceau \add{« admissible »} $\tilde{F}$ sur $\tilde{X}$, il faut avoir un triangle exact

\begin{tikzcd}[column sep=small]
 & Rf_!(\hat{\tilde{F}}) \arrow[dl] & \\
Rf_!(F) \arrow[rr] & & R\tilde{f}_*(\tilde{F}) \arrow[ul]
\end{tikzcd}

\page{162}
\note{p.~18 de l'auteur (crayon : 163).}
où $\hat{\tilde{F}}$ est le \emph{complété formel} de $\tilde{F}$ le long de $\tilde{X} - X$ (nb fini de points …). C'est ici, sauf erreur, que s'amorce le lien avec le corps de classes local, version cohomologique, dont je vais dire quelques mots maintenant.

E) \emph{Corps de classes local} comme formule de dualité.

Soit $V$ un anneau de val. discrète complet, à corps résiduel $k$ -- on suppose \struck{\ill{}} ou bien $k$ relevé en $k \subset V$ (donc $V \simeq k[[T]]$) -- ou bien $k$ parfait. Pour fixer les idées et être sur un terrain solide, je prends \add{d'abord} sur $K$ = corps des fractions de $V$ des coefficients \emph{finis} $F$ (comme sur $X$ tout à l'heure), et je m'intéresse aux $H^i(K, F)$ ou encore à $R\Gamma_K(F)$. L'essentiel du travail à faire consiste à définir une catégorie de coeff. convenable sur $k$ (sans doute la même que dans D)) et un foncteur
\[
F \longmapsto \mathbb{R}\underline{\Gamma}_K(F)
\]
vers ces coefficients, de telle façon que l'on ait
\[
\mathbb{R}\Gamma_K(F) \simeq R\Gamma_k(\mathbb{R}\underline{\Gamma}_K(F)) \,.
\]
\note{« ou bien $k$ parfait » : la version anglaise précise « of char.~$p>0$ », ajout de la copie dactylographiée. L'exposant de « $H^i(K, F)$ » est surchargé ; lu $i$. Dans la dernière formule les indices $K$ et $k$ sont mal distingués ; lus d'après le sens (et la version anglaise, lot~4, p.~65).}

\page{163}
\note{p.~19 de l'auteur (crayon : 164).}
Ceci correspond à l'intuition (sortant immédiatement des exemples élémentaires), moyennant quoi, pour $k$ alg. clos disons, les $H^0(K, F)$, $H^1(K, F)$ … ont une structure de groupes algébriques (ind-pro …) sur $k$. Dans cette construction, le schéma en anneaux sur $k$ défini par $V$ (introduit par Serre) et le « foncteur de Greenberg » (associant à un schéma sur $V$ un préschéma sur $k$) joueront un rôle essentiel.
\note{« le schéma en anneaux sur $k$ défini par $V$ » : ainsi sur la page ; la version anglaise (lot~4, p.~65) dit « the ring scheme of Witt vectors over $k$ ».}

Ceci posé, la formule de dualité s'écrira formellement comme dans (D) ci-dessus
\[
\Delta R\underline{\Gamma}_K(F) \simeq R\underline{\Gamma}_K(DF[1])
\]
où $D$ est la dualité de Cartier, $\Delta$ est la dualité de Serre. Pour un corps résiduel fini, elle devient (via « l'astuce de Lang » plus haut)
\[
\Delta_0 R\Gamma_K(F) \simeq R\Gamma_K(DF[2]) \,,
\]
$\Delta_0$ dualité de Pontrjagin. Elle contient encore le corps de classes local géométrique à la Serre, et le corps de classes local arithmétique sous sa forme classique.
\note{dans la seconde formule, les deux indices sont des $K$ ; la version anglaise dactylographiée porte $\Gamma_k$ à gauche. Le « 2 » du décalage est surchargé.}

\emph{Remarques} a) Si $F$ est premier à la car. résiduelle, la formule est très facile et bien connue. On peut la considérer comme un cas très particulier

\page{164}
\note{p.~20 de l'auteur (crayon : 165).}
de la « formule d'induction », pour un morphisme $i : s \to S$, en formalisme de dualité :
\[
i^!(D_S F) = D_s(i^*(F))
\]
(on prend ici l'inclusion de $s = \operatorname{Spec} k$ dans $S = \operatorname{Spec} V$). Donc le travail à faire concerne les coefficients $p$-primaires, $p = \operatorname{car} k > 0$. Le cas le plus subtil est celui d'inégales caractéristiques.

b) Le foncteur $R\underline{\Gamma}$ s'obtient en composant $Rj_*$ (où $j : U = \operatorname{Spec} K \to \operatorname{Spec} V = S$ est l'inclusion) avec une sorte de « foncteur de Greenberg », version cohomologique.

c) Dans (D) et (E), je me suis borné à parler de coefficients \emph{finis} $F$ -- c'est pour ceux-là que je suis sûr de mon fait. Mais il est certainement vrai que la formule de dualité est encore plus riche, qu'on peut encore énoncer quelque chose \add{p.~ex.} pour $F$ un schéma en groupes non nécessairement fini, p.~ex. un schéma abélien \add{(avec quelques fibres dégénérées dans (D) ?)} -- mais je n'ai jamais tiré cette question entièrement au clair, même heuristiquement. J'ai un vague souvenir d'un genre de formule
\note{dans la formule d'induction, l'indice du membre de gauche est un $S$ majuscule, celui de droite un $s$ minuscule.}

\page{165}
\note{p.~21 de l'auteur (crayon : 166).}
qui devrait être contenue dans le formalisme, savoir (si $k$ alg. clos disons) : si $F$ schéma abélien sur $K$, $F'$ le schéma abélien dual, et $G'$ le groupe pro-algébrique sur $k$ attaché « à la Greenberg » à son modèle de Néron, on a
\[
H^1(K, F) \overset{?}{\simeq} \mathrm{Ext}^1_{k\text{-gr}}(G', \mathbb{Q}/\mathbb{Z})
\]
(\emph{NB} sans aucune garantie). En principe, la \marginal{\uncertain{conj.} mentionnée sur les modèles de Néron devrait sortir de la machine de dualité locale}
\note{la fin de la phrase est écrite dans la marge gauche, de bas en haut, et rattachée par une accolade au texte ; le premier mot n'est lu qu'avec doute (la version anglaise, lot~4, p.~66 : « the previously mentioned duality conjecture concerning Neron models of SGA~7 »).}

d) Tu pourrais demander à Deligne s'il n'a pas approfondi les questions (D) et (E) dernièrement.
\note{l'alinéa d) est séparé de la suite par un trait horizontal.}

(F) \emph{Portée et limites de la topologie fppf}

Depuis les essais de Serre de trouver une « coh. de Weil » en utilisant la cohomologie d'un schéma à coefficients non seulement discrets $\mathbb{Z}/p^n\mathbb{Z}$ $(n \to \infty)$, ou $\mu_{p^n}$ $(n \to \infty)$, mais aussi continus (\add{p.~ex.} $W_n$, $n \to \infty$), qui donnaient de bons résultats pour récupérer un $H^1$ correct, pendant de nombreuses années j'étais \add{\uncertain{resté}} sous l'impression, que j'essayais en vain de préciser, qu'une « coh. de Weil » « $p$-adique » correcte, en car. $p > 0$, devrait pouvoir sortir d'une façon -- d'une autre de la cohomologie fppf, pour des coefficients finis p.~ex., ou des coefficients
\note{« $\mu_{p^n}$ » : le $\mu$ est mal formé. « d'une façon -- d'une autre » : un tiret tient la place de « ou ».}

\page{166}
\note{p.~22 de l'auteur (crayon : 167).}
plus généraux, groupes algébriques sur $k$. La construction en (B) du complexe des jacobiennes locales était, évidemment, \add{un peu} lié à cet espoir -- l'homologie donnerait peut-être ce que la cohomologie nous cèle ! Depuis lors, on dispose du formalisme de la cohomologie cristalline, et on sait que \add{(Berthelot)} (du moins \struck{\ill{}} pour $X$ projectif et lisse) elle a les propriétés correctes. Si on l'utilise comme une sorte d'étalon \uncertain{sûr} à quoi « mesurer » les autres cohomologies, on trouve ceci : la partie de la cohomologie cristalline $H^i_{\mathrm{cris}}(X)$ qui peut être décrite en termes de cohomologie fppf de $X$ à coefficients dans des $k$-groupes algébriques, n'est qu'une \struck{\ill{}} petite partie des $H^i$ ; de façon précise, utilisant la structure \add{très rigide} supplémentaire des $H^i$ (modules de t.f. sur l'anneau $W(k)$ des vecteurs de Witt) provenant de l'existence de \uncertain{l'isom} de Frobenius $H^i \xrightarrow{F} H^i$ (semi-linéaire), on trouve que l'on reste toujours dans la partie « de pente $\leqslant 1$ » (alors que les pentes possibles varient entre $0$ et $i$ …). Ceci explique pourquoi pour $i = 1$
\note{« (Berthelot) » est écrit au-dessus de « formalisme », avec un trait qui le renvoie après « on sait que » ; la place retenue est celle de la version anglaise.}

\page{167}
\note{p.~23 de l'auteur (crayon : 168).}
on a pu obtenir via fppf un $H^1$ correct, alors que pour le $H^2$ déjà tous les essais ont été infructueux.

À vrai dire, on conjecture que \emph{toute} la partie de pente $\leqslant 1$ des $H^i_{\mathrm{cris}}$ provient de fppf. Mais j'ai perdu complètement contact avec ces questions -- des gens, comme Mazur, Katz, Messing -- et bien sûr, Deligne -- devraient être au courant du point où en est la question.
\note{l'alinéa est suivi d'un triple trait horizontal.}

Ta question 7° semble indiquer qu'il y a malentendu en ton esprit sur la signification du « type d'homotopie » de $X$, $X$ un topos (p.~ex. le topos étale d'un schéma). Sans doute tu dois confondre l'algèbre homotopique qu'on peut faire sur $X$ à coups de faisceaux semi-simpliciaux, de champs en tous genres, et des relations entre ceux-ci -- et cet autre point de vue où $X$ (avec sa structure si riche de topos) disparaît virtuellement

\page{168}
\note{p.~24 de l'auteur (crayon : 169).}
pour ne devenir plus qu'un pâle élément d'une « catégorie homotopique » (ou pro-homotopique), déduite de celle des topos par un processus de localisation \ill{} -- à première vue, tout ce qui reste encore à ce pauvre $X$, ce sont ses $\pi_i$ -- et ses groupes de cohomologie à coefficients constants -- ou à la rigueur constants tordus. Quand on fouille un peu plus dans cette direction « que reste-t-il encore à ce pauvre $X$ » on tombe précisément sur les \add{$n$-}champs \emph{loc. constants} (car un $X \xrightarrow{f} X'$ qui est une équivalence d'homotopie induit une $(n+1)$-équivalence entre les cat. de $n$-champs loc. constants sur $X$ et sur $X'$) -- qui contiennent bien sûr les complexes \add{de chaînes} abéliens de longueur $n$ à faisceaux de coh. loc. constants, et l'hypercohomologie de ceux-ci. C'est ainsi qu'on arrive à \struck{\ill{}} \add{ce} triangle d'\struck{\ill{}} \add{\uncertain{objets}} se déterminant mutuellement
\note{après « localisation », deux ou trois mots (« à … ») ne sont pas lus ; la version anglaise dit « a very thorough process ». Le mot ajouté au-dessus du second biffé est d'une lecture très incertaine ; la version anglaise porte « this triangle of objects ».}

\begin{tikzcd}[column sep=tiny, row sep=2pt, nodes={font=\scriptsize}]
 & \text{topos (ou espace top., ou} & \\
 & \text{ens. ½simpl.) \emph{modulo}} & \\
 & \text{$n$-homotopie} \arrow[dddl, leftrightarrow] \arrow[dddr, leftrightarrow] & \\
 & & \\
 & & \\
\text{$n$-groupoïdes} & & \text{$n$-catégories « spéciales »} \\
\text{(à $n$-équiv. près)} \arrow[rr, leftrightarrow] & & \text{(à $n$-équivalence près)}
\end{tikzcd}

\note{« modulo » est souligné deux fois, « homotopie » une fois. Le triangle, dessiné à la main, est disposé ici ligne par ligne, comme dans la version anglaise (lot~4, p.~68).}

\page{169}
\note{p.~25 de l'auteur (crayon des archivistes : 170).}
On dit qu'une $n$-catégorie \add{$E_n$} est « \emph{spéciale} » (ou $n$-\emph{galoisienne}) si elle est $n$-équivalente à la cat. des $(n-1)$-champs loc. constants sur un espace top. (ou un topos) convenable, ou, ce qui doit revenir au même, si elle est $n$-équivalente à la catégorie des $n$-foncteurs $G_n \to (n\text{-Cat})$, où $G_n$ est un $n$-groupoïde. Si $X$, $G_n$, $E_n$ se correspondent ainsi, on dit que \struck{$E_n$} $G_n$ est le $n$-groupoïde fondamental de $X$, ou de $E_n$, ou que $E_n$ est la catégorie des $n$-systèmes locaux sur $X$, ou sur $G_n$, ou que $X$ est une réalisation géométrique de $G_n$ ou de $E_n$.
\note{« au même » est écrit « au $=$ » ; il note ainsi « même » tout au long du lot, rendu ici en toutes lettres. « $n$-systèmes locaux » : ainsi sur la page ; la version anglaise (lot~4, p.~68) corrige à la main en « $(n-1)$-systems ». Le symbole biffé avant $G_n$ est lu $E_n$ avec doute.}

En analogie avec le cas bien familier $n=1$, $G_n$ doit s'interpréter comme la sous-$n$-catégorie pleine de $\underline{\mathrm{Hom}}_n(E_n, (n\text{-cat}))$ formée des $n$-foncteurs $E_n \to (n\text{-cat})$ satisfaisant certaines propriétés d'exactitude (on a envie de dire : qui commutent aux $\varprojlim$ finies et aux $\varinjlim$ quelconques) ; mais cela nous attire dans l'inquiétant \uncertain{mirage} des $n$-limites dans les $n$-catégories (\emph{NB} le cas $n=2$ commence à nous être familier …). Il est prudent dans tout

\page{170}
\note{p.~26 de l'auteur (crayon : 171).}
ceci de supposer que $X$ est « localement homotopiquement trivial », ce qui assure que le pro-ensemble simplicial que lui associent Artin--Mazur (par le procédé des nerfs d'hyperrecouvrements) est essentiellement constant dans la catégorie homotopique ordinaire -- donc $X$ définit un type d'homotopie habituel. Ce n'est \uncertain{pas} le cas bien sûr pour le topos étale d'un schéma. Pour un tel cas, le $n$-groupoïde fondamental doit être conçu comme un \emph{pro-$n$-groupoïde} (rien d'étonnant à cela, vu la théorie familière des $\pi_1$) ; et $E_n$ comme une (ind)-$n$-catégorie (la ind-structure correspondant à l'exigence de locale trivialité, pour un $n$-champ variable, relativement à des recouvrements de plus en plus fins de $X$).
\note{« pro » est souligné deux fois, « $n$ » une fois. Le paragraphe est clos par un double trait.}

Je comprends cependant ta résistance instinctive à concevoir ce dépouillement extrême d'un beau topos $X$, au point de n'en retenir que le maigre type d'homotopie. Bien plus, je suis persuadé qu'en

\page{171}
\note{p.~27 de l'auteur (crayon : 172).}
allant jusqu'au bout de cette résistance, on arrive à généraliser et à approfondir la notion de « type d'homotopie », et à apporter encore une nouvelle eau au \ill{} développement d'un bon yoga homotopique. Voici où je veux en venir. \add{(\ill{})}
\note{la version anglaise porte ici « new grist to the mill » ; le mot avant « développement » n'est pas lu (il est trop court pour « moulin du »). Après « en venir », une insertion entre parenthèses, reliée par un trait à la ligne suivante, commence par « d' » et n'est pas lue ; elle n'a pas d'équivalent dans la version anglaise.}

Parlons d'abord de faisceaux, pour simplifier, au lieu de champs, \add{et \uncertain{plaçons-nous} sur le topos étale d'un schéma $X$.} Les faisceaux loc. constants -- modulo une condition de finitude supplémentaire assez anodine -- forment les plus simples des faisceaux \emph{constructibles}, à la définition desquels ils servent de modèle. Supposant $X$ cohérent (= qu.-qct et qu.-sép), les faisceaux constructibles généraux sont ceux pour lesquels il existe une partition \add{finie} $X = \bigcup_{i \in I} X_i$ de $X$ en « cellules » \add{ou « strates »} $X_i$, loc. fermées et constructibles, telles que \struck{\ill{}} la restriction de $F$ à chaque $X_i$ soit loc. constante \struck{et} (plus une condition de finitude …). Ainsi la catégorie des faisceaux constructibles sur $X$ (qui redonne la catégorie de tous les faisceaux en passant à une catégorie de ind-objets …) peut elle-même se concevoir comme limite inductive de catégories, associées à des partitions de plus en

\page{172}
\note{p.~28 de l'auteur (crayon : 173).}
plus fines de $X$. On peut donc, pour une telle partition \add{$\mathcal{P}$} fixée, se proposer une étude approfondie \add{de la cat.} des faisceaux (ou complexes de faisceaux, ou champs) qui sont « $\mathcal{P}$-constructibles » (ou plus généralement, qui sont « loc. constants » sur chacun des $X_i$). Ces catégories n'auront des structures vraiment satisfaisantes que si \struck{pour préciser \ill{}} elles sont stables par les opérations courantes -- telles que $R\underline{\mathrm{Hom}}$, ou $Rj_*j^*$ où $j : X_i \to X$ est une « cellule » de la partition\add{, etc}. En fait, si $X$ est excellent et qu'on dispose de la résolution des singularités, on sait que les faisceaux \add{de torsion} constructibles (sous réserve d'être premiers aux car. résiduelles) sont stables par \struck{\ill{}} ces opérations -- mais non avec une partition de $X$ fixée une fois pour toutes. Pour avoir une telle stabilité plus fine, il faut faire des hypothèses très strictes \struck{\ill{}} « d'équisingularité » de la stratification donnée de $X$, le long des strates. Je pense néanmoins qu'un affinement des techniques connues montrera que $X$ admet des stratifications arbitrairement fines ayant ces propriétés d'équisingularité (et avec des $X_i$ réguliers et connexes, mais peu importe \struck{\ill{}} cette propriété \uncertain{ici} …). À titre d'exemple
\note{« pour préciser » est biffé sous la ligne, à hauteur de « structures vraiment » (lecture incertaine) ; « de la cat. » est ajouté dans la marge gauche, avec un trait vers « des faisceaux ».}

\page{173}
\note{p.~29 de l'auteur (crayon : 174).}
supposons qu'il y ait exactement deux strates, une strate fermée $X_0$ et $X_1 = X - X_0$. D'après le dévissage d'Artin, un faisceau $F$ sur $X$ revient à la donnée d'un faisceau $F_0$ \add{$(= i_0^*(F))$} sur $X_0$, d'un faisceau $F_1$ \add{$(= i_1^*(F))$} sur $X_1$, et d'un homomorphisme $F_0 \to i_0^* i_{1*}(F_1)$. Pour que $F$ \add{($X_0 \overset{i_0}{\hookrightarrow} X \hookleftarrow X_1$)} soit $\mathcal{P}$-constructible, il faut et il suffit que $F_0$ et $F_1$ soient loc. constants (plus des conditions de finitude accessoires …), sur $X_0$ et $X_1$ respectivement. Alors (grâce aux hyp. d'équisingularité) il en sera de même de $i_0^* i_{1*}(F_1) = \varphi(F_1)$, et \add{la cat.} des faisceaux qui nous intéressent \struck{\ill{}} s'exprime entièrement en \struck{fonction} \add{termes} de catégories de faisceaux loc. constants sur $X_0$ et $X_1$, \add{i.e. des \uncertain{seuls}} \struck{\ill{}} \add{types} d'homotopie de $X_0$ et $X_1$, à cela près qu'il faut encore expliciter la nature du foncteur exact à g. $\varphi$. Je pense que ceci est possible, dans le contexte \add{schématique} où je me suis placé (techniquement assez sophistiqué), en introduisant un « voisinage tubulaire étale » de $X_0$ dans $X$ (qui est un topos fort intéressant, mais pas associé à un schéma). Mais cette construction technique ne fait que para-
\note{les deux exposants $i_0^*$, $i_1^*$ des ajouts sont mal formés et lus d'après la suite. Le diagramme des inclusions $X_0 \hookrightarrow X \hookleftarrow X_1$ est écrit dans la marge gauche, sous « Pour que $F$ ». Au milieu de la page, la suite « i.e. des seuls types d'homotopie » est faite d'ajouts interlinéaires sur un mot biffé et souligné ; l'ordre retenu est celui de la version anglaise (« i.e. of the mere homotopy type »). « exact à g. » : exact à gauche.}

\page{174}
\note{p.~30 de l'auteur (crayon : 175).}
phraser une intuition topologique extrêmement simple, que je vais expliciter, en supposant pour fixer les idées qu'on est sur le corps de base $\mathbb{C}$ et qu'on en profite pour travailler avec les espaces loc. compacts habituels. L'idée topologique derrière l'hypothèse d'équisingularité, c'est qu'il existe un \emph{voisinage tubulaire} $T$ de $X_0$ dans $X$ se rétractant sur $X_0$, et tel que $(X_0, T)$ au-dessus de $X_0$ soit un fibré localement trivial, donc $T - X_0$ au-dessus de $X_0$ est loc. trivial. En fait, si \add{$\partial T$} \struck{\ill{}} est le « bord » de $T$, qui sera aussi fibré loc. trivial sur $X_0$, $T$ sur $X$ est le fibré en cônes déduit $(\simeq (\partial T \times I) \amalg_{\partial T} X_0$, où $I = [0, 1]$, $\partial T \to \partial T \times I$ défini par $x \mapsto (x, 1)$, et $\partial T \to X_0$ est la projection) donc $\mathring{T} = T - X_0 \simeq \partial T \times [0, 1[$ est $X_0$-homotope à $\partial T$. Si $X_0$ et $X_1$ sont non singuliers, il en est de même de $\mathring{T}$ et $\partial T$, qui sont des fibrés topologiquement lisses sur $X_0$. D'ailleurs, \struck{\ill{}} posant $\tilde{X}_1 = X_1 - \mathring{T}$, \struck{\ill{}} l'inclusion $\tilde{X}_1 \hookrightarrow X_1$ est une équivalence d'homotopie, et $X$ se reconstitue, à homéomorphisme près, à partir du diagramme d'espaces

\begin{tikzcd}[column sep=large]
{(\mathring{T} = T - X_0 \simeq)\ \partial T} \arrow[r, "{j\ \text{(inclusion)}}"] \arrow[d, "{\text{fibration}\ p}"'] & {\tilde{X}_1\ (\simeq X_1)} \\
X_0 &
\end{tikzcd}

comme la somme amalgamée. En termes de ce diagramme d'espaces, le foncteur $\varphi$ plus
\note{« $T$ sur $X$ » : ainsi sur la page (et dans la version anglaise) ; on attendrait $T$ sur $X_0$. Le mot biffé après « si » est noirci ; « $\partial T$ » est écrit au-dessus.}

\page{175}
\note{p.~31 de l'auteur (crayon : 176).}
haut s'interprète immédiatement comme
\[
\varphi(F_1) \simeq p_* j^*(\tilde{F}_1)
\]
où $F_1 \mapsto \tilde{F}_1$ est la restriction de $X_1$ à $\tilde{X}_1$ (qui est donc une équivalence de catégories pour les faisceaux (loc. const.) envisagés). \struck{Il} \add{La donnée} d'\struck{\ill{}} $F = (F_0, F_1, u : F_0 \to \varphi F_1)$ peut donc aussi s'expliciter comme la donnée de
\[
F_0 \,,\ \tilde{F}_1 \,,\ \tilde{u} : p^*(F_0) \to j^*(\tilde{F}_1)
\]
où $F_0$ $(\tilde{F}_1)$ est un faisceau loc. constant sur $X_0$ (resp. $\tilde{X}_1$). Il faut se rappeler de plus qu'ici $p$ est une vraie fibration, et $j$ une inclusion (et pratiquement, -- \uncertain{en} $X_0$, $X_1$ lisses -- \struck{\ill{}} l'inclusion du bord dans une variété à bord).

Si on préfère, on peut aussi prendre le diagramme un peu moins joli (mais un peu plus canonique)

\begin{tikzcd}
\mathring{T} \arrow[r, "j'"] \arrow[d, "p'"'] & X_1 \\
X_0 &
\end{tikzcd}

revenant essentiellement au même, car formé d'espaces homotopes au précédent. On peut même y remplacer $X_0$ par $T$ ($X_0$ en étant un rétract par déformation) et écrire

\begin{tikzcd}
\mathring{T} \arrow[r, hook, "j'"] \arrow[d, hook, "p''"'] & X_1 \\
T &
\end{tikzcd}

\note{le diagramme est à cheval sur les pp.~31 et 32 de l'auteur : le sommet $T$ et la flèche $p''$ (un crochet d'inclusion) sont au bas de la p.~31, le reste en tête de la p.~32 ; il est donné ici en entier. « en $X_0$, $X_1$ lisses » : lecture incertaine de « en » (la version anglaise : « for the case »).}

\page{176}
\note{p.~32 de l'auteur (crayon : 177).}
où « pris à la lettre » $p''$ est maintenant une inclusion, mais « moralement » c'est \add{tjrs} une fibration à fibres très jolies [de dim. finie notamment, voire des variétés non singulières -- c'est donc beaucoup mieux que ce que nous donne le yoga de Cartan--Serre « toute application continue est équivalente à une fibration » …]. C'est ce dernier diagramme qui est susceptible d'une construction directe, purement algébrique, dans le cadre des schémas, lorsque l'on a développé la construction du voisinage tubulaire étale.

Le point où je veux en venir, c'est que la considération \add{de la cat.} des faisceaux \add{(ou complexes d'iceux, ou} \struck{(ou} \add{$n$-champs …)} $\mathcal{P}$-constructibles sur un $X$, où $\mathcal{P}$ est une stratification « équisingulière » donnée, revient dans le cas particulier actuel à la connaissance d'un diagramme de \emph{types d'homotopie} ordinaires \add{(ou de pro-types, si on revient à la top. étale)}

\begin{tikzcd}
\Delta \arrow[r, "j"] \arrow[d, "p"'] & X_1 \\
X_0 &
\end{tikzcd}

en prenant des systèmes de coeff locaux (ou de $n$-champs loc. constants) sur ces sommets \add{$X_0$, $X_1$} \struck{\ill{}}, reliés par des hom de compatibilité du type $p^*(F_0) \to j^*(F_1)$. Ce serait un exercice amusant (que je n'ai pas encore fait) de
\note{« tjrs » (toujours) est écrit au-dessus de « une fibration ». La parenthèse « (ou de pro-types … étale) » est écrite au-dessus de « ordinaires », qu'elle précède sur la page. Le sommet du diagramme est un $\Delta$, comme dans la version anglaise (lot~4, p.~73).}

\page{177}
\note{p.~33 de l'auteur (crayon : 178).}
vérifier et d'expliciter comment les « 6 opérations » habituelles sur les faisceaux (soit sur $X$, soit sur une réunion de strates de $X$) peuvent s'exprimer dans ce dictionnaire, dans le cas disons de strates non singulières (autrement, on aurait des difficultés avec les complexes dualisants, dont on aimerait qu'ils appartiennent à notre catégorie …) \add{et de revérifier le formulaire connu de ces opérations}. Mais il apparaît probablement que, pour bien pouvoir tout transcrire, il convient, plutôt \add{que} de regarder un diagramme du type

\begin{tikzcd}
\bullet \arrow[r] \arrow[d] & \bullet \\
\bullet &
\end{tikzcd}

dans la catégorie homotopique formée avec la cat. des ens. semi-simpliciaux, \add{de} regarder la catégorie des diagrammes d'ensembles semi-simpliciaux, et à partir de là passer à une catégorie homotopique de fractions.
\note{le diagramme « du type » est un simple coin de deux flèches, sans sommets écrits ; les points sont ajoutés ici. « semi-simpliciaux » est écrit la seconde fois « ½simpliciaux ».}

J'ai plus ou moins explicité dernièrement, en cogitant sur \uncertain{les} fondements de « topologie modérée » (i.e. où on élimine d'emblée tous les phénomènes de sauvagerie), comment une stratification équisingulière \struck{\ill{}}, à strates non singulières disons, d'un « espace modéré » \add{compact} donne naissance canoniquement à un diagramme \struck{\ill{}} d'espaces,
\note{un ajout interlinéaire au-dessus de « équisingulière » est biffé et cerclé, illisible.}

\page{178}
\note{p.~34 de l'auteur (crayon : 179).}
qui sont des variétés à bord, les \struck{\ill{}} flèches des diagrammes étant pour l'essentiel des fibrations loc. triviales de variétés à bord sur d'autres (à fibres des variétés à bord), \emph{et} des inclusions des bords dans les variétés à bord [en fait, on trouve des inclusions légèrement plus générales, certains des bords qui apparaissent étant munis d'une décomposition cellulaire « élémentaire », \add{i.e. \uncertain{dont}} \struck{\ill{}} les strates \add{fermées} sont encore des variétés à bord qui se recollent suivant des parties communes du bord ; et il y a lieu d'envisager \add{aussi} les inclusions de ces morceaux les uns dans les autres …], et se reconstitue à partir de ce diagramme \add{par recollement}. En d'autres termes, on a une description-dévissage canonique des espaces modérés compacts \add{$X$}, éventuellement munis de stratifications équisingulières à strates non sing., en termes de diagrammes \add{finis} de nature précise formés avec des variétés à bord. Quand on s'intéresse aux faisceaux \add{(ou complexes de faisceaux, ou $n$-champs …)} $\mathcal{P}$-constructibles sur $X$, où $\mathcal{P}$ est une telle stratification fixée, ceux-ci s'explicitent en termes du diagramme envisagé, dont il
\note{« élémentaire » : le mot porte une marque finale qui peut être un pluriel. Le mot biffé après l'ajout « i.e. dont » (lecture incertaine du second mot) est noirci.}

\page{179}
\note{p.~35 de l'auteur (crayon : 180).}
n'y a plus lieu de retenir que le « type d'homotopie ». On prévoit que les 6 opérations sur ces faisceaux peuvent se traduire ad hoc dans ce cadre homotopique. Enfin, si au lieu d'un seul espace modéré compact $X$, on a disons un morphisme modéré $f : X \to Y$ de tels, choisissant des stratifications équisingulières sur $X$ et $Y$ adaptées à $f$ (les strates de $X$ étant notamment des fibrés loc. triviaux sur celles de $Y$ …), on doit trouver un « morphisme » du diagramme de variétés à bord exprimant $X$ dans celui exprimant $Y$ (avec des morphismes mutuels se ramenant encore, pour l'essentiel, à des fibrations de variétés à bord sur d'autres) de telle façon que les 4 opérations $Rg_*$, $Rg_!$, $Lg^*$, $g^!$ entre faisceaux \add{$\mathcal{P}_X$- et $\mathcal{P}_Y$-constructibles} sur $X$ et $Y$ (ou sur des sous-espaces \add{$X'$, $Y'$} loc. fermés réunions de strates, de sorte que $f$ induise $g : X' \to Y'$) s'expriment encore en termes d'opérations standard entre seuls types d'homotopie. Enfin, toutes ces constructions, encore partiellement hypothétiques (y a du travail \add{de fondements} sur la planche !) doivent

\page{180}
\note{p.~36 de l'auteur (crayon : 181).}
pouvoir se paraphraser dans le cadre des schémas excellents, à grands coups de voisinages tubulaires étales. Dans l'un et l'autre cas, le « type d'homotopie fin » d'un espace modéré, resp. d'un schéma excellent, se définirait par passage à la limite à partir des \add{$\mathcal{P}$-}« types d'homotopie » associés à des stratifications équisingulières (à strates non sing.) $\mathcal{P}$ de plus en plus fines. Ce type d'homotopie \add{« fin »} incorporerait la connaissance, non des seuls faisceaux \add{ou $n$-}champs loc. constants, mais \add{(via un passage à la limite inductive)} la connaissance de \emph{tous}. Et il dépendrait, en un sens convenable, fonctoriellement de $X$. Dans le cas d'un schéma de type fini sur un corps alg. clos disons, les \emph{théorèmes de finitude cohomologiques} -- homotopiques les plus forts s'exprimeraient précisément en termes du type d'homotopie fin, et diraient que les types d'homotopie ordinaires qui en sont les éléments constituants sont essentiellement des « polyèdres finis » -- et même des variétés à bord compactes -- ou de façon plus précise,
\note{la phrase continue au lot~10.}

\end{document}
